\documentclass[12pt,a4paper]{article}
\usepackage[utf8]{inputenc}
\usepackage[spanish]{babel}
\usepackage{amsmath,amssymb}
\usepackage{geometry}
\usepackage{enumitem}
\usepackage{listings}
\usepackage{xcolor}
\usepackage{booktabs}

\geometry{margin=2cm}

\lstset{
  language=Python,
  basicstyle=\ttfamily\small,
  keywordstyle=\color{blue}\bfseries,
  stringstyle=\color{red},
  commentstyle=\color{gray}\itshape,
  showstringspaces=false,
  frame=single,
  breaklines=true,
  numbers=left,
  numberstyle=\tiny\color{gray}
}

\title{\textbf{Tecnología Digital VI --- Segundo Parcial} \\ Modelo de Examen 4}
\author{Universidad Torcuato Di Tella}
\date{}

\begin{document}
\maketitle
\noindent\textbf{Duración:} 2 horas \hfill \textbf{Nombre:} \hrulefill

\vspace{0.5cm}
\hrule
\vspace{0.5cm}

%% =============================================
\section*{Parte 1: Verdadero o Falso (2 pts c/u)}
\textit{Justifique brevemente las respuestas falsas.}

\begin{enumerate}[label=\alph*)]
    \item Una CNN tiene significativamente menos parámetros que un MLP para procesar imágenes, gracias al \textit{parameter sharing} y la localidad de los filtros.

    \item La fórmula del tamaño de salida de una convolución es $O = \frac{I - K + 2P}{S} + 1$. Si aplicamos un kernel de $5 \times 5$ con padding $P=2$ y stride $S=1$ a una imagen de $28 \times 28$, la salida tiene tamaño $28 \times 28$.

    \item Los Saliency Maps calculan la derivada de la predicción respecto a los pesos de la red para determinar qué píxeles son más importantes.

    \item En las GANs, el Generador y el Discriminador se entrenan cooperativamente para generar imágenes realistas.

    \item TF-IDF es una mejora sobre Bag-of-Words que penaliza palabras muy frecuentes en todo el corpus (como artículos y preposiciones), dando más peso a palabras discriminativas.

    \item En un Decoder-Only Transformer como GPT, se usa \textit{causal masking} para que cada token solo pueda atender a los tokens anteriores (y a sí mismo).
\end{enumerate}

\vspace{0.5cm}
\hrule
\vspace{0.5cm}

%% =============================================
\section*{Parte 2: Multiple Choice (3 pts c/u)}
\textit{Marque la opción correcta. Solo una es válida.}

\begin{enumerate}[label=\arabic*)]
    \item ¿Cuál es la principal contribución de la arquitectura \textbf{VGG}?
    \begin{enumerate}[label=\Alph*)]
        \item Demostró que usar siempre filtros de $3 \times 3$ con stride 1 es más eficiente que filtros grandes.
        \item Introdujo las skip connections para mitigar vanishing gradients.
        \item Fue la primera CNN en usar Dropout.
        \item Redujo los parámetros usando Global Average Pooling.
    \end{enumerate}

    \item ¿Qué es \textit{teacher forcing} en el entrenamiento de un modelo Seq2Seq?
    \begin{enumerate}[label=\Alph*)]
        \item Usar un modelo pre-entrenado (``teacher'') para inicializar los pesos del Seq2Seq.
        \item En cada paso del decoder, usar como input la palabra correcta del ground truth en vez de la predicción anterior.
        \item Forzar al encoder a producir un hidden state de mayor dimensión.
        \item Aplicar regularización L2 a los pesos del decoder.
    \end{enumerate}

    \item ¿Cuál de las siguientes analogías vectoriales es correcta usando embeddings de Word2Vec?
    \begin{enumerate}[label=\Alph*)]
        \item $v_{\text{madrid}} \approx v_{\text{paris}} - v_{\text{francia}} + v_{\text{españa}}$
        \item $v_{\text{madrid}} \approx v_{\text{paris}} + v_{\text{francia}} - v_{\text{españa}}$
        \item $v_{\text{madrid}} \approx v_{\text{paris}} + v_{\text{francia}} + v_{\text{españa}}$
        \item $v_{\text{madrid}} \approx v_{\text{españa}} - v_{\text{paris}} - v_{\text{francia}}$
    \end{enumerate}

    \item ¿Qué scheduler reduce el learning rate a la mitad cuando la loss de validación deja de mejorar durante varias épocas?
    \begin{enumerate}[label=\Alph*)]
        \item Step Decay
        \item Exponential Decay
        \item Cosine Annealing
        \item Reduce on Plateau
    \end{enumerate}

    \item En un Autoencoder, la \textbf{representación latente} es:
    \begin{enumerate}[label=\Alph*)]
        \item La salida final reconstruida del decoder.
        \item El vector comprimido que se obtiene entre el encoder y el decoder.
        \item Los pesos de la red entrenada.
        \item El vector de entrada después del preprocesamiento.
    \end{enumerate}
\end{enumerate}

\vspace{0.5cm}
\hrule
\vspace{0.5cm}

%% =============================================
\section*{Parte 3: Desarrollo (10 pts c/u)}

\begin{enumerate}[label=\arabic*)]
    \item \textbf{Convolución: cálculo paso a paso.} \\
    Dada una imagen de entrada de $6 \times 6$ (1 canal) y un kernel de $3 \times 3$:

    \[
    \text{Imagen} = \begin{pmatrix}
    1 & 0 & 1 & 0 & 1 & 0 \\
    0 & 1 & 0 & 1 & 0 & 1 \\
    1 & 0 & 1 & 0 & 1 & 0 \\
    0 & 1 & 0 & 1 & 0 & 1 \\
    1 & 0 & 1 & 0 & 1 & 0 \\
    0 & 1 & 0 & 1 & 0 & 1
    \end{pmatrix}
    \qquad
    \text{Kernel} = \begin{pmatrix}
    1 & 0 & -1 \\
    1 & 0 & -1 \\
    1 & 0 & -1
    \end{pmatrix}
    \]

    \begin{enumerate}[label=\alph*)]
        \item ¿Cuál es el tamaño de la salida si se usa padding $P=0$ y stride $S=1$?
        \item Calcule el valor de la posición $(0, 0)$ del feature map de salida (esquina superior izquierda).
        \item ¿Qué tipo de patrón detecta este kernel en particular?
    \end{enumerate}

    \item \textbf{Interpretabilidad: Grad-CAM vs Saliency Maps.} \\
    \begin{enumerate}[label=\alph*)]
        \item Explique cómo funcionan los Saliency Maps: ¿respecto a qué se calcula la derivada y qué se mantiene fijo?
        \item ¿Qué mejora aporta Grad-CAM respecto a Saliency Maps? ¿En qué capa de la red opera?
        \item Describa brevemente cómo funcionan los ataques adversarios: ¿cómo se pueden modificar los píxeles de una imagen para engañar a la red?
    \end{enumerate}

    \item \textbf{LSTMs y dependencias a largo plazo.} \\
    \begin{enumerate}[label=\alph*)]
        \item Explique por qué las RNNs simples sufren de vanishing gradients al entrenar con BPTT (Backpropagation Through Time).
        \item Describa la arquitectura interna de una celda LSTM: ¿qué son el cell state $C_t$ y el hidden state $h_t$? ¿Cuál es la función de cada gate (Forget, Input, Output)?
        \item ¿Por qué la operación de suma ``$+$'' en la actualización de $C_t$ (en lugar de multiplicación) es clave para mitigar el vanishing gradient? Haga una analogía con las skip connections de ResNet.
    \end{enumerate}
\end{enumerate}

\vspace{0.5cm}
\hrule
\vspace{0.5cm}

%% =============================================
\section*{Parte 4: Interpretación de Código (15 pts)}

Analice el siguiente código de PyTorch y responda las preguntas.

\begin{lstlisting}
import torch
import torch.nn as nn
import torchvision.models as models

# Cargar modelo pre-entrenado
resnet = models.resnet18(pretrained=True)

# Congelar todos los parametros
for param in resnet.parameters():
    param.requires_grad = False

# Reemplazar la ultima capa FC
num_features = resnet.fc.in_features
resnet.fc = nn.Sequential(
    nn.Linear(num_features, 256),
    nn.ReLU(),
    nn.Dropout(0.4),
    nn.Linear(256, 2)
)

criterio = nn.CrossEntropyLoss()
# Solo optimizar los parametros de la nueva capa fc
optimizador = torch.optim.Adam(resnet.fc.parameters(), lr=0.001)

for epoch in range(20):
    resnet.train()
    for imgs, labels in train_loader:
        optimizador.zero_grad()
        outputs = resnet(imgs)
        loss = criterio(outputs, labels)
        loss.backward()
        optimizador.step()
\end{lstlisting}

\begin{enumerate}[label=\alph*)]
    \item ¿Qué técnica de Deep Learning se está aplicando en este código? Descríbala brevemente.

    \item ¿Qué efecto tiene \texttt{param.requires\_grad = False} en el ciclo que recorre todos los parámetros? ¿Por qué se hace esto?

    \item ¿Por qué se pasa \texttt{resnet.fc.parameters()} al optimizador en vez de \texttt{resnet.parameters()}?

    \item La nueva capa \texttt{fc} tiene 2 neuronas de salida. ¿Cuántas clases tiene el problema? ¿Cuántas tenía el modelo original pre-entrenado en ImageNet?

    \item ¿Qué pasaría si las imágenes del \texttt{train\_loader} fueran en escala de grises (1 canal) en vez de RGB (3 canales)?

    \item Si el modelo estuviera sobreajustando al nuevo dataset, ¿qué estrategia adicional (además del Dropout ya incluido) podría implementar sin cambiar la arquitectura? Mencione al menos dos opciones.
\end{enumerate}

\end{document}
